% =====================================================================
\chapter*{Parte A. Cuestiones}
\phantomsection\addcontentsline{toc}{chapter}{Parte A. Cuestiones}
\markboth{Parte A. Cuestiones}{}
\label{cap:parteA}
% =====================================================================

Veinte preguntas de opción múltiple, de 1,5 puntos cada una (30 puntos en total). Cada pregunta tiene una única respuesta correcta; las respuestas erróneas no restan. La clave y su justificación están en el \cref{sec:clave}.

\begin{enumerate}[label=\textbf{\arabic*.},leftmargin=2em,itemsep=8pt,ref=\arabic*]

% ----------------------------------------------------------------- M07
\needspace{7\baselineskip}\item\label{q:1} \pts{1,5} Para $f(\theta)=2\theta^{2}$, el descenso de gradiente con tasa constante $\eta>0$ converge al mínimo desde cualquier $\theta_0$ si y solo si $\eta$ pertenece a
\begin{multicols}{4}
\begin{opciones}
  \item $(0;\ 0{,}25)$
  \item $(0;\ 0{,}5)$
  \item $(0;\ 1)$
  \item $(0;\ 2)$
\end{opciones}
\end{multicols}

\needspace{7\baselineskip}\item\label{q:2} \pts{1,5} En el método de la bola pesada, $v_{k+1}=\beta v_k+\nabla f(\theta_k)$ y $\theta_{k+1}=\theta_k-\eta v_{k+1}$, con $v_0=0$. Si $\beta=0{,}95$ y el gradiente es constante e igual a $g$, la velocidad $v_k$ tiende a
\begin{multicols}{4}
\begin{opciones}
  \item $0{,}95\,g$
  \item $1{,}95\,g$
  \item $20\,g$
  \item $95\,g$
\end{opciones}
\end{multicols}

\needspace{7\baselineskip}\item\label{q:3} \pts{1,5} Los ejemplos de un minibatch se muestrean de forma independiente y uniforme (con reemplazamiento). Si el tamaño del minibatch se multiplica por~4, la desviación típica de cada componente del gradiente estimado
\begin{multicols}{2}
\begin{opciones}
  \item se divide por~4.
  \item se divide por~2.
  \item no cambia.
  \item se multiplica por~2.
\end{opciones}
\end{multicols}

\needspace{7\baselineskip}\item\label{q:4} \pts{1,5} En AdamW con $\eta=10^{-3}$ y $\lambda=0{,}01$, ¿cuánto disminuye en un paso un peso $\theta=4$ por efecto del término de decaimiento de pesos?
\begin{multicols}{4}
\begin{opciones}
  \item $4\cdot 10^{-5}$
  \item $4\cdot 10^{-3}$
  \item $10^{-5}$
  \item Depende de $\hat v_t$.
\end{opciones}
\end{multicols}

\needspace{7\baselineskip}\item\label{q:5} \pts{1,5} El calentamiento lineal de la tasa de aprendizaje se usa sobre todo porque, al principio del entrenamiento,
\begin{opciones}
  \item los pesos son demasiado pequeños para que el gradiente sea informativo.
  \item los gradientes son grandes e irregulares y los momentos de Adam ($\hat m_t$, $\hat v_t$) se han estimado con muy pocos gradientes.
  \item la pérdida es convexa y admite una tasa mayor que al final.
  \item el dropout todavía no está activado.
\end{opciones}

% ----------------------------------------------------------------- M08
\needspace{7\baselineskip}\item\label{q:6} \pts{1,5} Si $X\T X$ es invertible, la solución de mínimos cuadrados de $\min_\theta\norm{X\theta-y}^{2}$ es
\begin{multicols}{2}
\begin{opciones}
  \item $\theta=X\T y$
  \item $\theta=(XX\T)^{-1}y$
  \item $\theta=(X\T X)^{-1}X\T y$
  \item $\theta=X^{-1}y$
\end{opciones}
\end{multicols}

\needspace{7\baselineskip}\item\label{q:7} \pts{1,5} En regresión logística, con $p=\sigma(w\T x+b)$ y pérdida $\ell=-\bigl[y\ln p+(1-y)\ln(1-p)\bigr]$, el gradiente $\nabla_w\ell$ es
\begin{multicols}{2}
\begin{opciones}
  \item $(p-y)\,x$
  \item $p(1-p)\,x$
  \item $(y-p)\,\sigma'(w\T x+b)$
  \item $x/p$
\end{opciones}
\end{multicols}

\needspace{7\baselineskip}\item\label{q:8} \pts{1,5} Un modelo de caracteres con $\abs{\V}=28$ que al inicio predice la distribución uniforme tiene una entropía cruzada inicial, en nats, de
\begin{multicols}{4}
\begin{opciones}
  \item $1$
  \item $2{,}8$
  \item $3{,}3322$
  \item $28$
\end{opciones}
\end{multicols}

\needspace{7\baselineskip}\item\label{q:9} \pts{1,5} Un clasificador obtiene una exactitud de $0{,}99$ en entrenamiento y de $0{,}62$ en validación. Lo más probable es que sufra
\begin{opciones}
  \item sesgo alto (infraajuste).
  \item varianza alta (sobreajuste).
  \item una tasa de aprendizaje demasiado baja.
  \item etiquetas erróneas en validación.
\end{opciones}

% ----------------------------------------------------------------- M09-M10
\needspace{7\baselineskip}\item\label{q:10} \pts{1,5} La capa \texttt{nn.Linear(30, 64)} de PyTorch implementa $z=xW+b$ con $x$ de tamaño $1\times 30$. ¿Qué forma tiene su atributo \texttt{weight}?
\begin{multicols}{4}
\begin{opciones}
  \item $(30, 64)$
  \item $(64, 30)$
  \item $(1, 64)$
  \item $(30, 30)$
\end{opciones}
\end{multicols}

\needspace{7\baselineskip}\item\label{q:11} \pts{1,5} La inicialización de Xavier con distribución normal para una capa $100\to 200$ usa una desviación típica de
\begin{multicols}{4}
\begin{opciones}
  \item $0{,}0577$
  \item $0{,}0816$
  \item $0{,}1$
  \item $0{,}1414$
\end{opciones}
\end{multicols}

\needspace{7\baselineskip}\item\label{q:12} \pts{1,5} ¿Cuál de estas activaciones toma un valor negativo en $t=-1$?
\begin{multicols}{4}
\begin{opciones}
  \item ReLU
  \item GELU
  \item Ambas.
  \item Ninguna.
\end{opciones}
\end{multicols}

\needspace{7\baselineskip}\item\label{q:13} \pts{1,5} La capa prealimentada SwiGLU de LLaMA~7B~\cite{llama} tiene $d=4096$, anchura oculta $11\,008$ y no tiene sesgos. ¿Cuántos parámetros tiene en cada bloque?
\begin{multicols}{4}
\begin{opciones}
  \item $90\,177\,536$
  \item $135\,266\,304$
  \item $67\,108\,864$
  \item $180\,355\,072$
\end{opciones}
\end{multicols}

\needspace{7\baselineskip}\item\label{q:14} \pts{1,5} Se quiere el jacobiano completo de $f\colon\R^{10}\to\R^{3}$ en modo inverso. ¿Cuántas pasadas hacia atrás hacen falta?
\begin{multicols}{4}
\begin{opciones}
  \item $1$
  \item $3$
  \item $10$
  \item $30$
\end{opciones}
\end{multicols}

\needspace{7\baselineskip}\item\label{q:15} \pts{1,5} Con números duales ($\varepsilon^{2}=0$), la parte en $\varepsilon$ de $(2+\varepsilon)^{3}$ vale
\begin{multicols}{4}
\begin{opciones}
  \item $8$
  \item $6$
  \item $12$
  \item $3$
\end{opciones}
\end{multicols}

% ----------------------------------------------------------------- M11
\needspace{7\baselineskip}\item\label{q:16} \pts{1,5} Con $\epsilon=0$ y ganancia $g=\uno$, para todo $x\ne 0$ el vector $\RMSNorm(10x)$, comparado con $\RMSNorm(x)$, es
\begin{multicols}{4}
\begin{opciones}
  \item 10 veces mayor.
  \item igual.
  \item 10 veces menor.
  \item Depende de $x$.
\end{opciones}
\end{multicols}

\needspace{7\baselineskip}\item\label{q:17} \pts{1,5} Con dropout invertido de probabilidad $p=0{,}25$, una activación $h=4$ vale durante el entrenamiento
\begin{multicols}{4}
\begin{opciones}
  \item $0$ o $4$
  \item $0$ o $5{,}3333$
  \item siempre $3$
  \item $0$ o $3$
\end{opciones}
\end{multicols}

\needspace{7\baselineskip}\item\label{q:18} \pts{1,5} El producto $260\cdot 260$, calculado en fp16 con redondeo al más próximo, da
\begin{multicols}{4}
\begin{opciones}
  \item $67\,600$
  \item $67\,584$
  \item \Inf
  \item $65\,504$
\end{opciones}
\end{multicols}

\needspace{7\baselineskip}\item\label{q:19} \pts{1,5} En fp16, $\exp(z)$ se desborda a infinito a partir de $z\approx$
\begin{multicols}{4}
\begin{opciones}
  \item $11{,}09$
  \item $88{,}72$
  \item $15$
  \item $65\,504$
\end{opciones}
\end{multicols}

\needspace{7\baselineskip}\item\label{q:20} \pts{1,5} Entrenando en fp32, la pérdida pasa de $2{,}1$ a \NaN\ de un paso al siguiente sin haber crecido antes. Según el árbol de diagnóstico de M11, §6.4, lo primero que debe buscarse es
\begin{opciones}
  \item una tasa de aprendizaje demasiado alta.
  \item un logaritmo, una división o una exponencial sin proteger.
  \item un desbordamiento de fp16.
  \item dropout activo durante la evaluación.
\end{opciones}

\end{enumerate}
